La implementación de una arquitectura de datos centralizada mejora significativamente la confiabilidad, disponibilidad y calidad de la información financiera, reduce la dependencia de procesos manuales y minimiza el riesgo de errores. El proyecto aporta beneficios como la automatización de procesos de integración de datos, la reducción del tiempo de generación de reportes financieros, el incremento en la confiabilidad de la información, la mejora en la toma de decisiones estratégicas, la optimización de los recursos del área financiera y la disponibilidad de información en tiempo real mediante dashboards interactivos. Asimismo, el uso de tecnologías cloud garantiza escalabilidad, disponibilidad y seguridad en la infraestructura tecnológica.
